\chapter{Происхождение пространственного электромагнитного ядра}
\label{ch:v8-baryon-em-spatial-kernel-origin-gate}

Зарядовое тождество предыдущей главы не содержит координат. Поэтому
проверим, какие части нейтрон-протонного знака защищены одной зарядовой
алгеброй, а какие требуют отдельного пространственного состояния.

\section{Трёхточечный электростатический оператор}

Пусть $g_{rs}=g_{sr}$ --- среднее пространственного ядра между копиями
$r$ и $s$. На фиксированной нумерации копий общий оператор имеет вид
\begin{equation}
 H_{\rm el}=\frac1T\left[
 \mu\sum_{r=1}^{3}Q_r^2
 +2\sum_{1\le r<s\le3}g_{rs}Q_rQ_s
 \right].
 \label{eq:v8-baryon-em-spatial-general-operator}
\end{equation}
Коэффициент $\mu$ включает локальную часть, а $g_{rs}$ --- матричный
элемент ядра Грина в трёхчастичном пространственном состоянии. Эти числа
не определяются спектром самого оператора заряда.

При соглашении $p=(u,u,d)$ и $n=(u,d,d)$ точный расчёт даёт
\begin{equation}
 T(E_n-E_p)
 =-\frac{\mu}{3}+\frac{2}{3}(g_{23}-2g_{12}).
 \label{eq:v8-baryon-em-labelled-spatial-splitting}
\end{equation}
Следовательно, даже положительные $g_{rs}$ без перестановочной симметрии
не защищают знак. Например,
\begin{equation}
 \mu=1,\qquad g_{12}=g_{13}=1,\qquad g_{23}=3
 \quad\Longrightarrow\quad
 T(E_n-E_p)=\frac13>0.
 \label{eq:v8-baryon-em-positive-kernel-sign-counterexample}
\end{equation}

\section{Что восстанавливает перестановочное усреднение}

Усредним протонное состояние по трём положениям заряда $-1/3$, а
нейтронное --- по трём положениям заряда $2/3$. Тогда для каждой пары
\begin{equation}
 \left\langle Q_rQ_s\right\rangle_p=0,
 \qquad
 \left\langle Q_rQ_s\right\rangle_n=-\frac19.
 \label{eq:v8-baryon-em-pair-charge-permutation-average}
\end{equation}
Обозначая
\begin{equation}
 \bar g=\frac{g_{12}+g_{13}+g_{23}}{3},
 \label{eq:v8-baryon-em-spatial-kernel-average}
\end{equation}
получаем
\begin{equation}
 T(E_n-E_p)=-\frac{\mu+2\bar g}{3}.
 \label{eq:v8-baryon-em-symmetrized-spatial-splitting}
\end{equation}
Поэтому отрицательный электростатический знак действительно устойчив при
$\mu>0$, $\bar g>0$, но его защита использует перестановочно-усреднённое
пространственно-зарядовое состояние, а не только тождество полного заряда.

\begin{proposition}[Условная защита электростатического знака]
В перестановочно-усреднённом трёхкварковом состоянии любое симметричное
положительное парное ядро входит в нейтрон-протонную разность только через
$\bar g$ и даёт $E_n-E_p<0$ при $T,\mu,\bar g>0$.
\end{proposition}

\section{Масштабная орбита пространственного состояния}

Даже если выбрать кулоновское ядро $G(r)=1/|r|$, его матричный элемент не
становится безразмерной константой. Для нормированной волновой функции в
$d$ пространственных измерениях положим
\begin{equation}
 \psi_s(r_1,r_2,r_3)=s^{3d/2}\psi(sr_1,sr_2,sr_3),
 \qquad s>0.
 \label{eq:v8-baryon-em-wavefunction-dilation}
\end{equation}
Тогда точно
\begin{equation}
 \left\langle\frac1{|r_i-r_j|}\right\rangle_{\psi_s}
 =s\left\langle\frac1{|r_i-r_j|}\right\rangle_{\psi}.
 \label{eq:v8-baryon-em-coulomb-dilation}
\end{equation}
Без кинетического и связывающего пространственного родителя параметр $s$
не выбирается. Компактное пространство прежних томов также не устраняет
пробел: в проекте нет отображения трёх кварковых копий в три координатных
оператора и нет сцепления их зарядов со скалярным ядром Грина.

\begin{nogocriterion}[Пространственное ядро не следует из зарядов]
Наличие операторов $Q_r$, общего следа и эпсилон-синглета не определяет
$g_{rs}$, радиальный масштаб или отношение локальной и парной частей.
Пространственный член в предыдущем гейте отсутствовал вследствие
усечения модели, а не вследствие доказанного равенства нулю.
\end{nogocriterion}

\section{Итог гейта}

\begin{protocol}
\begin{enumerate}
 \item зарядовая арифметика: \textbf{точная};
 \item знак без перестановочного усреднения: \textbf{не защищён};
 \item положительная контрмодель смены знака: \textbf{построена};
 \item знак после перестановочного усреднения: \textbf{отрицателен условно};
 \item величина $\bar g$: \textbf{не выведена};
 \item кулоновская масштабная орбита: \textbf{доказана точно};
 \item координатный трёхкварковый носитель в проекте: \textbf{не найден};
 \item универсальная пространственная электромагнитная энергия:
       \textbf{не получена};
 \item следующий гейт: \textbf{происхождение магнитного сверхтонкого члена}.
\end{enumerate}
\end{protocol}

Машинный сертификат сохранён в
\path{s2t_v8_baryon_em_spatial_kernel_origin_gate_results.json}.